\documentclass[a4paper, 12pt]{article}

\usepackage{utils}

\renewcommand*{\today}{27 February 2026}

\begin{document}

\hotbox{Analyse 4}{CM 8}{\today}

\subsection{Rayon de convergence}

\begin{definition}
    On définit $D_r \{z \in \mathbb{C}, |z| < r\}$ avec $r \in \R_+^*$.
\end{definition}

\begin{theoreme}{}{}
    Soit $r \in \R_+^*$.
    S'il existe $M > 0$ tel que
    $$
    \forall n \in \N, |a_n| r^n \leq M
    $$
    Alors pour tout $z \in D_r$, la série entière $\sum_{k = k_0}^{+\infty} a_k z^k$ converge absolument.
\end{theoreme}

\begin{demonstration}
    Pour tout $n \in \N$ on a
    $$
    |a_n z^n| = |a_n| |z^n| = |a_n| r^n \left|\frac{z}{r}\right|^n \leq M \frac{|z|^n}{r^n}
    $$
    donc $\sum\limits_{k \geq k_0} |a_k z^k| \leq M \sum\limits_{k \geq k_0} \left|\frac{z}{r}\right|^k$ une série géométrique de raison $\dfrac{|z|}{r} < 1$.
\end{demonstration}

\begin{definition}
    On définit le \textbf{rayon de convergence} d'une série entière $\sum\limits_{k = k_0}^{+\infty} a_k z^k$ comme étant
    $$
    R = \sup_{r > 0} \{ r \in \R_+^* | \sum\limits_{k = k_0}^{+\infty} a_k z^k \text{ converge pour tout } z \in D_r\}
    $$
\end{definition}

\begin{remarque}
    Ainsi on a
    \begin{itemize}
        \item si $|z| < R$ la série entière converge absolument,
        \item si $|z| = R$ la série entière peut converger ou diverger,
        \item si $|z| > R$ la série entière diverge.
    \end{itemize}
\end{remarque}

\begin{remarque}
    Si a partir d'un certain rang $n$ on a $a_n = 0$, alors $R = +\infty$.
\end{remarque}

\begin{theoreme}{}{}
    Soit $R$ le rayon de convergence de la série entière $\sum\limits_{k = k_0}^{+\infty} a_k z^k$. Alors
    $$
    \frac{1}{R} = \limsup_{n \to +\infty} \sqrt[n]{|a_n|} = \inf_{n_0 \in \N} \sup_{n \geq n_0} \sqrt[n]{|a_n|}
    $$
\end{theoreme}

\begin{demonstration}
    \begin{description}
        \item[Si $r < R$.] 
        La série de terme général $|a_n|r^n$ est convergente,
        à partir d'un certain rang $n_0$ on a
        $$
        \forall n > n_0, |a_n| r^n \leq 1 \iff \sqrt[n]{|a_n|} \leq \frac{1}{r}
        $$
        On peut passer au sup, d'où $\sup_{n \geq n_0} \sqrt[n]{|a_n|} \leq \frac{1}{r}$
        puis à la limite $n_0 \to +\infty$ d'où $\limsup_{n \to +\infty} \sqrt[n]{|a_n|} \leq \frac{1}{r}$.
    
        \item[Si $r > R$.]
        On fait pareil

        \item[Si $r = R$.]
        On ne peut rien dire.
    \end{description}
    Ainsi avec $c_n = a_n b^n$ on a
    $$
    \sqrt[n]{|c_n|} = \sqrt[n]{|a_n|} r < 1 \iff \sqrt[n]{|a_n|} < \frac{1}{r}
    $$
\end{demonstration}

\begin{corollaire}{}{}
    Soit $(a_n)_{n \in \N}$ une suite de nombres réels ou complexes telle que $\dfrac{|a_{n+1}|}{|a_n|}$ converge.

    Alors la série entière $\sum\limits_{k = k_0}^{+\infty} a_k z^k$ admet un rayon de convergence $R$ donné par
    $$
    \frac{1}{R} = \lim_{n \to +\infty} \frac{|a_{n+1}|}{|a_n|}.
    $$
\end{corollaire}

\begin{demonstration}
    On pose $c_n = a_n b^n$.
    Alors
    \begin{align*}
        \dfrac{|c_{n+1}|}{|c_n|} = \frac{|a_{n+1}|}{|a_n|} r &\iff \frac{|a_{n+1}|}{|a_n|} < \frac{1}{r} \\
        &\implies \lim_{n \to +\infty} \frac{|a_{n+1}|}{|a_n|} \leq \frac{1}{r}
    \end{align*}
\end{demonstration}

\begin{proposition}{}{}
    Soit $R$ le rayon de convergence de la série entière $\sum\limits_{k = k_0}^{+\infty} a_k z^k$ tel que $R > 0$.
    Soit $r \in ]0, R[$.

    Alors on a
    $$
    \forall \varepsilon > 0, \exists N \in \N, \forall n > N, \forall z \in D_r, \left| \sum_{k = n_0}^n a_k z^k - \sum_{k = n_0}^{+\infty} a_k z^k \right| < \varepsilon
    $$
    C'est à dire que le reste de la série entière converge uniformément sur $D_r$.
    et donc que $z \mapsto \sum\limits_{k = k_0}^{+\infty} a_k z^k$ est continue sur $D_r$.
\end{proposition}

\begin{demonstration}
    Soit $r' \in ]r, R[$ tel que $|z| < r'$.
    Pour tout $n > N$ on a
    $$
    |a_n z^n| \leq |a_n| r^n \leq |a_n| r'^n \left(\frac{r}{r'}\right)^n 
    $$
    Pour tout $z \in D_r$ on a
    \begin{align*}
        \left| \sum_{k = n_0}^n a_k z^k - \sum_{k = n_0}^{+\infty} a_k z^k \right| &= \sum_{k = n+1}^{+\infty} |a_k z^k| \\
        &\leq \sum_{k = n+1}^{+\infty} |a_k| |z|^k \\
        &\leq \sum_{k = n+1}^{+\infty} M \left(\dfrac{r}{r'}\right)^k \\
        &\leq M \sum_{k = n+1}^{+\infty} \left(\dfrac{r}{r'}\right)^k \\
        &\leq M \dfrac{r^{n+1}}{r'^{n+1}} \dfrac{1}{1 - \frac{r}{r'}} \\
        &\leq \dfrac{Mr'}{r' - r} \left(\dfrac{r}{r'}\right)^{n+1} \quad \text{qui ne dépend plus de z}
    \end{align*}
    Donc $\lim_{n \to +\infty} \left(\dfrac{r}{r'}\right)^{n+1} = 0$ d'où le résultat.
\end{demonstration}

\begin{proposition}{}{}
    Soit $\sum\limits_{k = k_0}^{+\infty} a_k z^k$ et $\sum\limits_{k = k_0}^{+\infty} b_k z^k$ deux séries entières de rayons de convergence respectifs $R_a$ et $R_b$.

    Pour tout $\forall (\alpha, \beta) \in \mathbb{C}^2$ et pour tout $z \in D_{\min(R_a, R_b)}$ on a
    \begin{enumerate}
        \item 
        $$
        \alpha \sum_{k = k_0}^{+\infty} a_k z^k + \beta \sum_{k = k_0}^{+\infty} b_k z^k = \sum_{k = k_0}^{+\infty} (\alpha a_k + \beta b_k) z^k
        $$

        \item
        $$
        \left(\sum_{k = k_0}^{+\infty} a_k z^k\right) \left(\sum_{k = k_0}^{+\infty} b_k z^k\right) = \sum_{k = k_0}^{+\infty} \left( \sum_{i = k_0}^k a_i b_{k - i} \right) z^k
        $$
    \end{enumerate}
\end{proposition}

\begin{demonstration}
    \begin{description}
        \item[1.]
        Simple.

        \item[2.]
        Soit $z \in D_{\min(R_a, R_b)}$.
        On a
        \begin{align*}  
            \sum_{i = i_0}^n |c_i| |z|^i = \sum_{i = i_0}^n \left| \sum_{h + k = i} a_k b_h \right| |z|^i \\
            &= \sum_{i = i_0}^n \left| \sum_{h + k = i} a_k b_h z^{k + h}\right| \\
            &= \sum_{i = i_0}^n \left| \sum_{h + k = i} |a_k| |z^k| |b_h| |z^h|\right| \\
            &\leq \left( \sum_{k = k_0}^n |a_k| |z|^k \right) \left( \sum_{h = k_0}^n |b_h| |z|^h \right) \\
            &\leq \left( \sum_{k = k_0}^{+\infty} |a_k| |z|^k \right) \left( \sum_{h = k_0}^{+\infty} |b_h| |z|^h \right)
        \end{align*}
        Les sommes partielles $\sum\limits_{k = k_0}^n c_k z^k$ sont uniformément bornées et elles sont croissantes donc
        $\left(\sum\limits_{k = k_0}^n |c_k| |z^k|\right)_n$ est absolument convergente, son reste tend vers 0
        et donc $\left(\sum\limits_{k = k_0}^{+\infty} c_k z^k\right)_n$ converge absolument. 
    \end{description}
\end{demonstration}

\begin{remarque}
    Le rayon de convergence de la nouvelle série entière peut être plus grand que $\min(R_a, R_b)$.
\end{remarque}

\end{document}